% =============================================================================
%  Part V: Numerical Implementation  (Sections 23--26)
% =============================================================================

\part*{Part V: Numerical Implementation}

\section{Hybrid Computation Strategy}

The volume integrals derived in Parts~I and~II can be evaluated analytically
by SymPy (symbolic computer algebra) for both the sphere and cube geometries.
However, for deployment in a numerical wavefield-modelling code, a pure-NumPy
implementation is desirable.  This section describes the hybrid numerical
strategy used in the \texttt{CubicTMatrix} Python package.

\subsection{The Singularity Structure of the Integrands}

The two volume integrals that define the local site T-matrix have
fundamentally different singularity structures at the origin:

\textbf{Green's tensor integral} (Eq~\ref{eq:Gamma0} / Section~13.1): The integrand
$G_{ij}(\bx) \sim 1/r$ as $r \to 0$. In three dimensions,
$d^3x \sim r^2\,dr$, so the integral converges absolutely. Standard 3D
Gauss--Legendre quadrature evaluates it accurately.

\textbf{Second-derivative integral} (Eq~\ref{eq:ABCcube}): The integrand
$G_{ij,kl}(\bx) \sim 1/r^3$ as $r \to 0$. This is the static
Eshelby singularity. The volume integral \emph{diverges}---it must be
understood as a Cauchy principal value plus a shape-dependent depolarisation
tensor. Direct numerical quadrature captures the divergent static part with
an $n$-dependent finite remainder, producing incorrect results.

\subsection{The Hybrid Approach}

The solution is to use a different method for each integral:
\begin{itemize}
\item \textbf{$\Gamma_0$: full numerical quadrature.}  Vectorised 3D
  Gauss--Legendre quadrature over $[-a,a]^3$ with $n_{\mathrm{GL}} = 32$
  points per dimension.  The GL nodes never include the origin (they are
  the zeros of Legendre polynomials on $(-1,1)$), and the $1/r$
  singularity is integrable, so the integral converges rapidly.  This
  correctly captures both the real (static) and imaginary
  (frequency-dependent) contributions.

\item \textbf{$A^c$, $B^c$, $C^c$: polynomial Taylor expansion.}  We
  decompose the Green's tensor into a singular static part ($1/r^3$)
  and a smooth polynomial remainder.  Only the smooth remainder is
  integrated---the static singularity drops out automatically because it
  is excluded from the Taylor series of the frequency-dependent terms.
  The polynomial integrand is evaluated by GL quadrature, which
  integrates it exactly (to Taylor truncation order).
\end{itemize}
This hybrid strategy achieves the accuracy of the SymPy symbolic
computation (matching $T_1$ and $T_2$ to $\sim 10^{-8}$, $T_3$ to
$\sim 0.3\%$) while running in pure NumPy at negligible computational
cost.


%=====================================================================
\section{Smooth Green's Tensor Decomposition}
%=====================================================================

\subsection{Splitting off the Singularity}

The Green's tensor (Eq~\ref{eq:Grad}) can be decomposed as:
\begin{equation}
  G_{ij}(\bx)
  = G_{ij}^{\mathrm{static}}(\bx)
  + G_{ij}^{\mathrm{s}}(\bx),
\end{equation}
where the static part $G_{ij}^{\mathrm{static}} \sim 1/r$ produces the
divergent $1/r^3$ second derivatives, and the smooth part
$G_{ij}^{\mathrm{s}}$ is an entire function of $\bx$ (a convergent
power series).

The smooth part arises from the odd-power terms in the Taylor expansion of
$f(r)\cdot r$ and $g(r)\cdot r$ about $r=0$.  It has the structure:
\begin{equation}
  \boxed{%
    G_{ij}^{\mathrm{s}}(\bx)
    = \delta_{ij}\,\Phi(r^2) + x_i\,x_j\,\Psi(r^2)
  }
\end{equation}
where $\Phi(u) = \sum_n \phi_n\,u^n$ and $\Psi(u) = \sum_n \psi_n\,u^n$
are power series in $u = r^2$.

\subsection{Taylor Coefficients}

The Taylor coefficients are derived by expanding the exponentials
$e^{i\omega r/\alpha}$ and $e^{i\omega r/\beta}$ in the definitions of
$f(r)$ and $g(r)$ (Eqs~\ref{eq:fdef}--\ref{eq:gdef}), multiplying by $r$, and collecting terms of
odd power in $r$ (which become polynomial after division by $r$):
\begin{equation}
  \boxed{%
    \phi_n = \frac{(i\omega/\beta)^{2n+1}}{(2n+1)!\;4\pi\rho\beta^2}
  }
\end{equation}
\begin{equation}
  \boxed{%
    \psi_n = \frac{1}{4\pi\rho}
    \left[
      \frac{(i\omega/\alpha)^{2n+3}}{\alpha^2}
      - \frac{(i\omega/\beta)^{2n+3}}{\beta^2}
    \right]
    \frac{1}{(2n+3)!}
  }
\end{equation}
In practice, $N_{\text{Taylor}} = 8$ terms suffice for convergence at the
frequencies of interest ($\omega a/\beta < 1$).

\subsection{Smooth Second Derivatives}

The second spatial derivatives of $G_{ij}^{\mathrm{s}}$ are computed via the
chain rule, using $\partial r^2/\partial x_k = 2x_k$.  The three independent
components needed for the cube's $A$, $B$, $C$ decomposition (Eq~\ref{eq:ABCdef}) are:
\begin{align}
  G_{11,11}^{\mathrm{s}}
    &= 4x_1^2\,\Phi'' + 2\Phi' + 2\Psi + 10x_1^2\,\Psi'
       + 4x_1^4\,\Psi'',
  \\
  G_{11,22}^{\mathrm{s}}
    &= 4x_2^2\,\Phi'' + 2\Phi' + 2x_1^2\,\Psi'
       + 4x_1^2 x_2^2\,\Psi'',
  \\
  G_{12,12}^{\mathrm{s}}
    &= \Psi + 2(x_1^2 + x_2^2)\,\Psi'
       + 4x_1^2 x_2^2\,\Psi'',
\end{align}
where primes denote $d/du$ with $u = r^2$.  These are polynomial in
$(x_1, x_2, x_3)$---completely free of singularities at the
origin---and are integrated over the cube by GL quadrature.

The integrals give $A^c$, $B^c$, $C^c$ via (Eq~\ref{eq:ABCdef}):
\begin{equation}
  A^c = \int G_{11,22}^{\mathrm{s}}\,d^3x, \quad
  B^c = \int G_{12,12}^{\mathrm{s}}\,d^3x, \quad
  C^c = \int G_{11,11}^{\mathrm{s}}\,d^3x - A^c - 2B^c.
\end{equation}


%=====================================================================
\section{Displacement-Traction Basis}
%=====================================================================

For integration into the Kennett reflectivity propagator framework
(Section~20.2), the T-matrix must operate on the displacement-traction
state vector rather than the strain tensor.  This section derives the
$6\times 6$ T-matrix in the
$(u_z, u_x, u_y, t_{zz}, t_{xz}, t_{yz})$ basis.

\subsection{Coordinate System}

We adopt a right-handed coordinate system with $z$~(down), $x$~(right),
$y$~(out of page).  The mapping from the generic $(1,2,3)$ Cartesian indices
used in Parts~I--IV to $(z,x,y)$ is: $1\to z$, $2\to x$, $3\to y$.  The
Voigt ordering is:
\begin{equation}
  \boldsymbol{\varepsilon}_V
  = [\varepsilon_{zz},\;\varepsilon_{xx},\;\varepsilon_{yy},\;
     2\varepsilon_{xy},\;2\varepsilon_{zy},\;2\varepsilon_{zx}]^T.
\end{equation}

\subsection{Strain Extraction from Displacement-Traction}

For a plane wave with horizontal wavenumbers $(k_x,k_y)$, the six
independent strains can be expressed in terms of the six
displacement-traction components \emph{without requiring~$k_z$}.  The
traction components absorb the vertical-wavenumber dependence through the
stress-strain relation.  The key relations are:
\begin{align}
  \varepsilon_{zz}
    &= \frac{t_{zz} - \lambda\,ik_x\,u_x - \lambda\,ik_y\,u_y}
            {\lambda + 2\mu},
  \\
  \varepsilon_{xx} &= ik_x\,u_x, \qquad
  \varepsilon_{yy} = ik_y\,u_y,
  \\
  2\varepsilon_{zx} &= \frac{t_{xz}}{\mu}, \qquad
  2\varepsilon_{zy} = \frac{t_{yz}}{\mu},
  \\
  2\varepsilon_{xy} &= ik_y\,u_x + ik_x\,u_y.
\end{align}
These define a $6\times 6$ strain extraction matrix $\bS$ such that
$\boldsymbol{\varepsilon}_V = \bS(k_x,k_y)\;\bv^0$
where $\bv = (u_z, u_x, u_y, t_{zz}, t_{xz}, t_{yz})^T$:
\begin{equation}
\bS
= \begin{pmatrix}
  0 & -\dfrac{\lambda\,ik_x}{M} & -\dfrac{\lambda\,ik_y}{M}
    & \dfrac{1}{M} & 0 & 0 \\[8pt]
  0 & ik_x & 0 & 0 & 0 & 0 \\
  0 & 0 & ik_y & 0 & 0 & 0 \\
  0 & ik_y & ik_x & 0 & 0 & 0 \\
  0 & 0 & 0 & 0 & 0 & \dfrac{1}{\mu} \\[8pt]
  0 & 0 & 0 & 0 & \dfrac{1}{\mu} & 0
\end{pmatrix},
\end{equation}
where $M = \lambda + 2\mu$.  Note that $u_z$ does not appear explicitly:
the vertical strain $\varepsilon_{zz}$ is extracted entirely from the
traction $t_{zz}$ and the horizontal displacement components.

\subsection{Voigt T-Matrix}

The $6\times 6$ Voigt T-matrix (Eqs~\ref{eq:TVblock}--\ref{eq:Sblock}) in the $(z,x,y)$ coordinate
system is:
\begin{equation}
  \bT_V
  = \begin{pmatrix} \bD & \mathbf{0} \\
                     \mathbf{0} & \bS_{\mathrm{shear}} \end{pmatrix},
  \qquad
  \bD = T_1\,\bJ + (2T_2 + T_3)\,\bI_3,
  \qquad
  \bS_{\mathrm{shear}} = 2T_2\,\bI_3.
\end{equation}
This maps Voigt strain $\to$ Voigt effective stress perturbation:
$\Delta\boldsymbol{\sigma}^*_V
= \bT_V\,\boldsymbol{\varepsilon}^0_V$.

\subsection{Effective Stiffness Contrast Matrix}

The effective stiffness contrast tensor $\Delta c^*_{mnlp}$ in Voigt form
has cubic symmetry:
\begin{equation}
  \Delta\mathbf{c}^*_V
  = \begin{pmatrix} \bD^* & \mathbf{0} \\
                     \mathbf{0} & \bS^* \end{pmatrix},
\end{equation}
where
\begin{equation}
  D^*_{ii} = \Delta\lambda^* + 2\Delta\mu^{*,\mathrm{diag}}, \quad
  D^*_{ij} = \Delta\lambda^* \; (i \neq j), \quad
  \bS^* = 2\Delta\mu^{*,\mathrm{off}}\,\bI_3,
\end{equation}
with the effective contrasts $\Delta\lambda^*$,
$\Delta\mu^{*,\mathrm{diag}}$, $\Delta\mu^{*,\mathrm{off}}$ given by
Eqs~\eqref{eq:Dmuoff}--\eqref{eq:Dlamstar_c}.

\subsection{Traction Projection}

The traction on a $z = \text{const}$ surface is extracted from Voigt stress
by a $3\times 6$ projection matrix~$\bP$:
\begin{equation}
  \begin{pmatrix} t_{zz} \\ t_{xz} \\ t_{yz} \end{pmatrix}
  = \bP\,\boldsymbol{\sigma}_V, \qquad
  \bP
  = \begin{pmatrix}
    1 & 0 & 0 & 0 & 0 & 0 \\
    0 & 0 & 0 & 0 & 0 & 1 \\
    0 & 0 & 0 & 0 & 1 & 0
  \end{pmatrix},
\end{equation}
where $\sigma_{zx}$ is Voigt index~6 and $\sigma_{zy}$ is Voigt index~5.

\subsection{The $6\times 6$ T-Matrix in the $(\bu,\bt)$ Basis}

The full T-matrix in the displacement-traction basis maps
$(\bu,\bt)^0 \to$ scattered $(\bu,\bt)$.  Its
structure is:
\begin{equation}
  \boxed{%
    \bT_{6\times 6}
    = V\,\bigl[\text{density block} + \text{stiffness block}\bigr]
  }
\end{equation}
The \textbf{density scattering} contribution to the scattered displacement
is:
\begin{equation}
  u_i^{\mathrm{scat}}
  = V\,\omega^2\,\Delta\rho^*\,\Gamma_0\;u_i^0,
\end{equation}
giving a diagonal $3\times 3$ block in the upper-left of
$\bT_{6\times 6}$.

The \textbf{stiffness scattering} contribution to the scattered traction
is:
\begin{equation}
  \bt^{\mathrm{scat}}
  = V\;\bP\;\Delta\mathbf{c}^*_V\;\bS(k_x,k_y)\;
    \bv^0.
\end{equation}
This fills the lower $3\times 6$ block.  The complete matrix is:
\begin{equation}
  \bT_{6\times 6}
  = V
  \begin{pmatrix}
    \omega^2\,\Delta\rho^*\,\Gamma_0\,\bI_3 & \mathbf{0}_3 \\
    \multicolumn{2}{c}{%
      \bP\;\Delta\mathbf{c}^*_V\;\bS(k_x,k_y)
    }
  \end{pmatrix}.
\end{equation}
The displacement scattering from stiffness contrasts (upper-right blocks)
requires the strain-Green's tensor $H_{ijk}$ evaluated at the observation
point, which is handled by the propagation framework (FFTProp).  At the
local site, the self-consistent amplification factors (Eqs~\ref{eq:Aeoff}--\ref{eq:Aucube}) already
account for this coupling.


%=====================================================================
\section{Numerical Example (Cube --- NumPy)}
%=====================================================================

Using the same parameters as Section~10:
$V_P = 5000$~m/s, $V_S = 3000$~m/s, $\rho = 2500$~kg/m$^3$.
Contrasts: $\Delta\lambda = 2$~GPa, $\Delta\mu = 1$~GPa,
$\Delta\rho = 100$~kg/m$^3$.
Frequency $f = 10$~Hz, cube half-width $a = 10$~m.

\medskip
\begin{center}
\begin{tabular}{ll}
\toprule
\textbf{Quantity} & \textbf{Value (NumPy hybrid method)} \\
\midrule
$k_P a$ & $0.1257$ \\
$k_S a$ & $0.2094$ \\
\midrule
$\Gamma_0$ & $2.604567\mathrm{e}{-09} + 4.347360\mathrm{e}{-10}\,i$ \\
$A^c$      & $0.000000\mathrm{e}{+00} - 8.607635\mathrm{e}{-14}\,i$ \\
$B^c$      & $0.000000\mathrm{e}{+00} + 3.973958\mathrm{e}{-14}\,i$ \\
$C^c$      & $0.000000\mathrm{e}{+00} - 1.044478\mathrm{e}{-19}\,i$ \\
\midrule
$T_1^c$    & $0.000000\mathrm{e}{+00} + 2.252429\mathrm{e}{-04}\,i$ \\
$T_2^c$    & $0.000000\mathrm{e}{+00} - 4.633677\mathrm{e}{-05}\,i$ \\
$T_3^c$    & $0.000000\mathrm{e}{+00} - 2.088956\mathrm{e}{-10}\,i$ \\
\midrule
$A_u$                    & $1.001029\mathrm{e}{+00} + 1.719804\mathrm{e}{-04}\,i$ \\
$A_\theta$               & $9.999997\mathrm{e}{-01} + 5.830547\mathrm{e}{-04}\,i$ \\
$A_e^{\mathrm{off}}$     & $1.000000\mathrm{e}{+00} - 9.267354\mathrm{e}{-05}\,i$ \\
$A_e^{\mathrm{diag}}$    & $1.000000\mathrm{e}{+00} - 9.267375\mathrm{e}{-05}\,i$ \\
\midrule
$\Delta\rho^{*,c}$           & $1.001029\mathrm{e}{+02} + 1.719804\mathrm{e}{-02}\,i$ \\
$\Delta\lambda^{*,c}$        & $1.999999\mathrm{e}{+09} + 1.616595\mathrm{e}{+06}\,i$ \\
$\Delta\mu^{*,\mathrm{diag}}$ & $1.000000\mathrm{e}{+09} - 9.267375\mathrm{e}{+04}\,i$ \\
$\Delta\mu^{*,\mathrm{off}}$  & $1.000000\mathrm{e}{+09} - 9.267354\mathrm{e}{+04}\,i$ \\
\midrule
$\Delta\mu^{*,\mathrm{diag}} - \Delta\mu^{*,\mathrm{off}}$
  & $-3.886223\mathrm{e}{-05} - 2.088956\mathrm{e}{-01}\,i$ \\
\bottomrule
\end{tabular}
\end{center}

\medskip
The values of $T_1^c$ and $T_2^c$ agree with the SymPy symbolic computation
to approximately $10^{-8}$ relative error.  $T_3^c$ and $C^c$ differ by
approximately $0.3\%$, because the NumPy implementation retains higher-order
Taylor terms (8~terms \textit{vs.}\ the $O((\omega a)^7)$ truncation in the
symbolic code).

\bigskip
\noindent
\textit{Generated from \texttt{CubicTMatrix/effective\_contrasts.py} and
\texttt{CubicTMatrix/voigt\_tmatrix.py} using NumPy hybrid computation.
$\Gamma_0$ evaluated by 3D Gauss--Legendre quadrature ($32^3$~points);
$A^c$, $B^c$, $C^c$ evaluated by polynomial Taylor expansion (8~terms)
with GL quadrature.}
